\section{Related Work}
\label{sec:related}

\textbf{Static Model Merging.} Static model merging combines multiple independently fine-tuned models into a single set of weights without additional training~\cite{wortsman2022model, jin2022dataless}. Prominent methods include \textit{Task Arithmetic}~\cite{ilharco2022editing}, which directly adds task vectors to a base model, and \textit{TIES-Merging}~\cite{yadav2023ties}, which resolves parameter sign conflicts and redundant values by masking and scaling. Other works, such as DARE~\cite{yu2023language} and RegMean~\cite{choshen2022fusing}, utilize randomized pruning or activation-covariance matching. While static merging techniques are elegant and require zero parameters at inference, they are fundamentally static: they produce a single, compromise model that cannot dynamically adjust its weights on a sample-by-sample basis.

\textbf{Dynamic Routing and Mixture of Experts.} To support input-dependent adaptation, Mixture of Experts (MoE) architectures employ a routing function to route token or sample representations to specialized sub-networks~\cite{shazeer2017outrageously, fedus2022switch}. However, standard MoE routing functions contain trainable routing layers (e.g., linear projections with Softmax gates) that must be trained jointly with the model parameters using sophisticated auxiliary load-balancing losses~\cite{zoph2022designer, megablocks}. In the post-hoc ensembling paradigm, recent works like QWS-Merge~\cite{gong2024qmerge} attempt to train parametric dynamic ensembling routers using wave-superposition theory. However, as we demonstrate, these trainable gating networks suffer from severe small-sample inductive overfitting when trained on small calibration datasets, collapsing under vectorized online deployment.

\textbf{Parameter-Free Routing and Orthogonalization Limits.} To perform dynamic ensembling without parametric layers, we propose \textbf{Parameter-Free Task-Space Projection (PFSR)}. PFSR eliminates trainable parameters by projecting representation features directly onto raw task centroids derived from the expert classifiers. However, because independently trained specialists often share overlapping feature or representation subspaces, their raw centroids are rarely orthogonal, which has led researchers to speculate that raw centroids cause severe \textit{routing cross-talk}. A natural, theoretically elegant candidate to resolve this cross-talk is L{\"o}wdin-Orthogonalized Task-Space Projection (OTSP), which computes a perfectly decoupled, orthonormal task coordinate system. However, in this work, we deconstruct this approach to show that under standard isotropic noise, OTSP's orthogonalization is either mathematically redundant or actually degrades performance due to noise amplification. This establishes PFSR as the optimal, most robust parameter-free dynamic router.

\textbf{L{\"o}wdin Symmetric Orthogonalization.} Originally proposed in quantum chemistry to orthogonalize overlapping atomic wavefunctions~\cite{lowdin1950non}, L{\"o}wdin orthogonalization is a symmetric, order-invariant orthogonalization technique. Unlike Gram-Schmidt orthogonalization, which is highly order-dependent (the first vector is unaltered while subsequent vectors are projected away), L{\"o}wdin orthogonalization treats all input vectors symmetrically~\cite{soderstrom1974theory}. It solves the constrained optimization problem of finding an orthonormal basis that minimizes the squared Euclidean distance to the original non-orthogonal vectors. Despite its mathematical elegance, its application to representation-space task projection and parameter-free model ensembling remains entirely unexplored.
